\documentclass{article}
\usepackage[left=2cm, right=2cm, top=2cm, bottom=2cm]{geometry}
\usepackage{graphicx}
\usepackage{amsmath}
\usepackage{amssymb}
\usepackage{amsthm}
\usepackage{fancyhdr}
\usepackage{verbatim}
\usepackage{listings}
\usepackage{xcolor}
\usepackage{pdfpages}
\usepackage{cancel}
\usepackage{physics}

\lstset{
    language=C++,
    basicstyle=\ttfamily\footnotesize,
    keywordstyle=\color{blue},
    commentstyle=\color{green},
    stringstyle=\color{red},
    numbers=left,
    numberstyle=\tiny\color{gray},
    frame=single,
    breaklines=true,
    captionpos=b,
}


\title{MTH 553: Lecture \# 23}
\author{Cliff Sun}

\newtheorem{theorem}{Theorem}[section]
\newtheorem{lemma}[theorem]{Lemma}
\newtheorem{definition}[theorem]{Definition}
\newtheorem{conjecture}[theorem]{Conjecture}
\newtheorem{proposition}[theorem]{Proposition}
\newtheorem{corollary}[theorem]{Corollary}
\newtheorem{one minute paper}[theorem]{One Minute Paper}

\pagestyle{fancy}
\lhead{\textbf{\thepage}\ \ \nouppercase{\rightmark}}
\chead{MTH 553: Lecture \# 23}
\rhead{Cliff Sun}

\begin{document}

\maketitle

\section*{Lecture Span}
\begin{enumerate}
    \item Cauchy problem for 2nd order
\end{enumerate}

\section*{Cauchy problem for 2nd order}

\begin{align*}
    u_{xx} + u_{yy} &= u \\
    u(x,0) &= f(x) \\
    u_y(x,0) &= g(x)
\end{align*}

Then try to find a power series solution. Note, if $u$ is real analytic (i.e. it equals its Taylor series) around some point $(\tilde{x},0)$ on the real axis, then $u$ is 

\begin{equation}
    u(x,y) = \sum_{j,k=0}^{\infty}\left( \partial_x^{j}\partial_{y}^{k}u \right)(\tilde{x},0)\frac{1}{j!k!}(x - \tilde{x})^{j}y^k
\end{equation}

Note, this is just the Taylor series. Hence 

\begin{equation}
    f(x) = u(x,0) = \sum_{j=0}^{\infty}\left( \partial_{x}^{j}u \right)(\tilde{x},0)\frac{1}{j!}(x - \tilde{x})^{j}
\end{equation}

So $f$ is \underline{real analytic} and similarly so is $g$. \textbf{Q:} is real analyticity of Cauchy data $f,g$ sufficient for getting a solution? \textbf{A:} yes!

\begin{theorem}
    \underline{Cauchy-Kowalewski Theorem}. In essence, "Analytical Cauchy Problem has unique local analytic solution". Assume $f$ and $g$ are real analytic, then consider the following:
    \begin{equation}
        u_{yy} = G(x,y,u,u_x,u_y,u_{xx},u_{yx})
    \end{equation}
    In essence, you can rewrite the differential equation for $u_{yy}$. Moreover, 
    \begin{align}
        u(x,0) &= f(x) \\
        u_y(x,0) &= g(x)
    \end{align}
    Where $G$ is real analytic as a function of $7$ variables near $(\tilde{x},0, f(\tilde{x}), f'(\tilde{x}), g(\tilde{x}), f''(\tilde{x}), g'(\tilde{x}))$.
    Then, a unique real analytic solution $u(x,y)$ exists in a neighborhood of $(\tilde{x},0)$. 
\end{theorem}

\begin{proof}
    Evans section 4.6. Rough idea: assume $u$ exists and determine all derivatives. Then determine that the power (Taylor) series converges. 
\end{proof}

Note, this solution does not guarantee \underline{well-posedness}. 

\begin{definition}
    A PDE is \underline{well-posed} if a solution
    \begin{enumerate}
        \item exists
        \item is unique 
        \item depends continuously on the data
    \end{enumerate}
    For instance, slightly perturbing $f$ and $g$ will slightly perturb $u(x,y)$. 
\end{definition}

Usually (often), if a PDE problem is ill-posed, then it is a bad model of physical reality. This is not always the case. 

\end{document}